% LaTeX fragment. All proofs identified as ours are written for this run.
% Suggested wrapper: article; amsmath, amssymb, amsthm, mathtools, hyperref.
% Bibliography locators are also recorded in REFERENCES_CHECKED.md.
\providecommand{\ord}{\operatorname{ord}}
\providecommand{\Tr}{\operatorname{Tr}}
\providecommand{\aco}{\operatorname{aco}}
\providecommand{\N}{\mathbb N}
\providecommand{\C}{\mathbb C}

\section*{Conventions and dependency discipline (ours)}
All scalar entire functions are on $\C$. Unless expressly stated otherwise they are nonzero; zeros are counted with algebraic multiplicity. For a zero at the origin of multiplicity $m$ we first divide by $z^m$. Constants and nonzero polynomials have order zero; when an algebra is discussed the zero function is included by convention, but no genus is assigned to it. Write
\[
 E_p(w)=(1-w)\exp\left(\sum_{j=1}^p\frac{w^j}{j}\right),\quad p\in\N\cup\{0\},
 \qquad M_f(r)=\max_{|z|=r}|f(z)|.
\]
Empty products are one, empty zero families have exponent zero, and the infimum of an empty set of positive exponents is $+\infty$. In logarithmic formulae for order we may replace $\log\log M$ by $\log\log\max(e,M)$; this only repairs the constant/small-radius convention. Genus means the manuscript's maximum of the minimal canonical-product degree and the degree of the residual polynomial, not just the degree of the primary factors. Statements about general vector-valued sequences that the page does not define are labelled reconstructions; they are not attributed to the author.

The elementary complex-analysis tools used below are Cauchy's integral formula, the mean-value property for harmonic functions, the maximum principle, and Liouville's theorem. For completeness the particular Jensen formula needed throughout is derived in H0DT.11. The exact published Fredholm zero-multiplicity theorem used in H0DT.1 is Grothendieck's Bourbaki expos\'e 91, \S2, Theorem 2, p.~381 (talk March 1954; scanned text revised June 1957); the Hilbert-space product is independently stated in \S4(c), p.~383. This is not an invocation of Lidskii's trace identity to prove itself.

\section*{H0DT.1 --- The Hilbert-space determinant and its genus}
\textbf{Precise setting.} $H$ is a complex Hilbert space and $A$ is trace class, with singular values $s_j$ and nonzero eigenvalues $\lambda_j$, the latter repeated by algebraic multiplicity. Set $d_A(z)=\sum_{n\ge0}\alpha_n(A)z^n$, where $\alpha_n(A)=\Tr(\Lambda^n A)$ and $\alpha_0=1$. The statement is Known: Grothendieck, expos\'e 91, \S4(c), p.~383, and Koutsonikos-Kouloumpis and Lesch (2024), \S1.2, p.~2, (1.1) and the following paragraph, state exactly this product and genus zero. The separability assumed in the latter source causes no restriction: the closed span of the ranges of $A$ and $A^*$ is separable, invariant for both, and $A$ is zero on its orthogonal complement.

\textbf{Derivation, ours, with the stated spectral theorem as input.} A singular-value decomposition gives singular values $s_{i_1}\cdots s_{i_n}$ for $\Lambda^n A$. Hence
\[
 |\alpha_n(A)|\le \|\Lambda^n A\|_1
 =\sum_{i_1<\cdots<i_n}s_{i_1}\cdots s_{i_n}
 \le \frac{\|A\|_1^n}{n!}.
\]
The trace inequality used here follows by writing a trace-class operator as $\sum_j s_j v_j\otimes u_j^*$ and applying $|\langle v_j,u_j\rangle|\le1$. These bounds prove local uniform convergence of $d_A$ and coefficient domination by $h(z)=\prod_j(1+s_jz)$.

For any finite collection of nonzero eigenvalues, take the finite-dimensional invariant sum of their generalized eigenspaces. Schur triangularization in an orthonormal basis $e_1,\ldots,e_N$ makes the selected eigenvalues the diagonal entries $\langle Ae_i,e_i\rangle$. Choose diagonal phases $\eta_i$ so that $\eta_i\langle Ae_i,e_i\rangle=|\lambda_i|$. The finite-rank contraction $D$ having these phases on that basis and zero on the perpendicular space satisfies
\[
 \sum_{i=1}^N|\lambda_i|=|\Tr(DA)|\le\|A\|_1.
\]
Increasing the finite collection proves $\sum_j|\lambda_j|\le\|A\|_1$. Grothendieck's zero-multiplicity theorem identifies the zeros of $d_A$ as $-1/\lambda_j$, with the stated multiplicities. H0DT.21, proved below without the determinant formula, now gives genus zero. Since $d_A(0)=1$,
\[
 d_A(z)=\prod_j(1+z\lambda_j).
\]
Absolute local uniform convergence permits comparison of Taylor coefficients and yields
\[
 \alpha_n(A)=\sum_{i_1<\cdots<i_n}\lambda_{i_1}\cdots\lambda_{i_n},\qquad
 \Tr A=\sum_j\lambda_j.
\]
The identity for $|A|$ gives the singular-value coefficient sum written on the page. All infinite sums here converge absolutely. This proves the theorem and the equivalence of the product and all coefficient identities. The trace identity alone is only the $n=1$ consequence; it is not by itself equivalent without applying it to all exterior powers.

\section*{H0DT.2 --- Coefficient and radial determinant majorants}
\textbf{Proof, ours.} The singular-value calculation in H0DT.1 proves the entire chain $|\alpha_n(A)|\le\alpha_n(|A|)=\|\Lambda^n A\|_1\le\|A\|_1^n/n!$ for every $n\ge0$. This is the live Corollary 3 on scan 71 and the coefficient assertion repeated on 79. Consequently
\[
 |d_A(z)|\le\sum_n|\alpha_n(A)||z|^n
 \le\sum_n\alpha_n(|A|)|z|^n=d_{|A|}(|z|).
\]
If the page uses $h=(AA^*)^{1/2}$ instead, its nonzero eigenvalues are the same singular values, so its determinant is the same. This comparison is radial: it does not assert $|d_A(z)|\le|d_{|A|}(z)|$ at the same complex argument.

\section*{H0DT.3 --- The cancelled Hilbert-space product}
The cancellation is documentary, not a mathematical refutation. Under the Hilbert trace-class hypotheses actually written, H0DT.1 proves both $\det(I+A)=\prod_j(1+\lambda_j)$ and the entire-parameter identity. Their coefficient and trace consequences follow by coefficient comparison there. The same exact assertions appear in Grothendieck, expos\'e 91, \S4(c), p.~383. No assertion is made here for every nuclear operator on every Banach space; that would change the hypotheses.

\section*{H0DT.4 --- The cancelled reduction to the scalar majorant}
The proposed reduction is valid: absolute eigenvalue summability, the radial coefficient majorant, and the scalar proposition H0DT.21 give H0DT.1. The margin explicitly announces the stronger absolute-sum estimate elsewhere; it must not be deleted from the reading.

There is also a direct way to make the projected-trace estimate suffice. For every finite eigenvalue subfamily, projection onto its finite-dimensional invariant generalized eigenspace gives
\[
 \left|\sum_{j\in F}\lambda_j\right|\le\|A\|_1.
\]
Crucially, this is a bound for arbitrary finite subfamilies, not just the first $N$ values of one enumeration. If a complex family satisfies $|\sum_{j\in F}\lambda_j|\le C$ for every finite $F$, split any finite family according to the signs of its real parts: both positive and negative real-part sums have magnitude at most $C$, so $\sum|\Re\lambda_j|\le2C$. Likewise $\sum|\Im\lambda_j|\le2C$. Therefore
\[
 \sum_{j\in F}|\lambda_j|\le4C
\]
for every finite $F$, proving absolute summability. The phase contraction in H0DT.1 improves the constant to one. For algebraic multiplicities, use invariant root spaces and Schur triangularization; eigenvectors alone account only for geometric multiplicities. This is the explicitly supplied spectral completion of the sketch.

The scalar family $(-1)^j/j$ has bounded ordered partial sums but divergent absolute sum. It does \emph{not} have uniformly bounded sums over all finite subfamilies (take positive terms). It is consequently not a counterexample to the finite-subfamily inference available here. An earlier draft of our audit overlooked that distinction; the final verdict is True as written for the reduction, with this correction documented in the audit trail. For positive $h$, the direct elementary-symmetric spectral expansion supplies the product, completing the last step without assuming the general determinant identity.

\section*{H0DT.5 --- Weierstrass factorization}
\textbf{Proof, ours.} The zeros $a_n\ne0$ of a nonzero entire function with $f(0)\ne0$ have no finite accumulation point. For $|w|\le1/2$ choose the logarithm at $w=0$. Then
\[
 \log E_p(w)=-\sum_{j=p+1}^\infty\frac{w^j}{j},\qquad
 |\log E_p(w)|\le2|w|^{p+1}.
\]
Choose $p_n=n$. For each $R$, eventually $|a_n|\ge2R$, so $\sum_n(R/|a_n|)^{n+1}<\infty$. The logarithmic series therefore converges absolutely uniformly on $|z|\le R$ after omission of finitely many factors. Its exponential is nonzero; multiplication by the finitely many omitted factors shows that $P(z)=\prod_nE_{p_n}(z/a_n)$ is entire with exactly the prescribed zeros and multiplicities. The quotient $q=f/P$, with removable singularities filled in, is entire and nowhere zero. Since $\C$ is simply connected, $q'/q$ has an entire primitive $g_0$. The derivative of $q e^{-g_0}$ is zero, so a choice of logarithm of its nonzero constant value gives $q=e^g$. This proves the factorization. The empty/finite zero family is treated by a finite product and the same quotient argument. Uniform absolute convergence of the tail product is the intended meaning even on a compact set containing a zero; it does not mean an absolutely convergent logarithm at the zero itself.

\section*{H0DT.6 --- Admissible degrees and a constant degree}
\textbf{Proof, ours.} If nonnegative integers $p_n$ satisfy $\sum_n(R/|a_n|)^{p_n+1}<\infty$ for every $R>0$, the logarithmic estimate in H0DT.5 proves convergence with those degrees. The quotient argument again supplies $g$. In particular $\sum_n|a_n|^{-k-1}<\infty$ permits the constant choice $p_n=k$, and it permits the least such nonnegative integer $p\le k$. The page mixes $p$ in the preceding sentence and $k$ in its display: these describe two admissible choices, with correspondingly different functions $g$. They need not give the same $g$.

\section*{H0DT.7 --- Genus is well defined}
Fix the least nonnegative integer $p$ for which $\sum_n|a_n|^{-p-1}<\infty$. The product with these fixed primary factors is independent of enumeration by absolute local uniform convergence. If $e^{g_1}P=e^{g_2}P$, then $e^{g_1-g_2}=1$. The continuous entire function $g_1-g_2$ takes values in the discrete set $2\pi i\mathbb Z$, hence is constant. Thus polynomiality and degree of $g$ do not depend on the logarithm chosen. If $g$ is a polynomial, $G=\max(p,\deg g)$ is a well-defined integer; take constants to have degree zero for this purpose. If no finite $p$ exists, or $g$ is not a polynomial, the function has no finite genus. A factor $z^m$ does not alter this definition. This proves the claims attached to the definition, with the inserted word \emph{grand} retained. The uncorrected deleted word \emph{petit} would give, for $e^z$, zero instead of genus one, which explains why the visible correction matters.

\section*{H0DT.8 --- Increasing the common primary degree}
For $r\ge p$,
\[
 E_r(w)=E_p(w)\exp\left(\sum_{j=p+1}^r\frac{w^j}{j}\right).
\]
Each $\sum_n a_n^{-j}$ in this finite range converges absolutely. Thus
\[
 \prod_nE_r(z/a_n)=\prod_nE_p(z/a_n)
 \exp\left(\sum_{j=p+1}^r\frac{z^j}{j}\sum_na_n^{-j}\right).
\]
If $G\le r$, subtract the displayed polynomial from $g$; its degree is at most $r$. This gives the representation claimed on scan 108. Conversely a fixed-degree $r$ canonical representation with residual polynomial of degree at most $r$ gives genus at most $r$ after decreasing the primary degree by the same identity. This is our explicit resolution of the doubtful terminal indices, not a new reading of them.

\section*{H0DT.9 --- Genus zero}
With $f(0)\ne0$, genus zero means $p=0$ and constant $g$, hence exactly
\[
 f(z)=C\prod_n(1-z/a_n),\qquad C\ne0,\quad\sum_n|a_n|^{-1}<\infty.
\]
The converse follows immediately from H0DT.7. Without the standing nonvanishing assumption the complete statement is $f(z)=Cz^m\prod_n(1-z/a_n)$. The polynomial $f(z)=z$ refutes the version without that factor if the assumption is removed. This extension is ours and is used explicitly where a later page allows a zero at the origin.

\section*{H0DT.10 --- Order and eventual growth}
For nonconstant entire $f$, the maximum modulus tends to infinity, and
\[
 \rho=\limsup_{r\to\infty}\frac{\log\log M_f(r)}{\log r}\in[0,\infty].
\]
If $k>\rho$, then eventually $\log\log M_f(r)<k\log r$, so $M_f(r)<e^{r^k}$. Conversely such an eventual bound for some finite $k$ gives $\rho\le k$. Taking infima over admissible positive $k$ recovers $\rho$, including order zero. Nonnegativity follows from any nonzero Taylor coefficient of positive degree $n$: Cauchy's estimate gives $M_f(r)\ge |c_n|r^n$, and $\log(n\log r+O(1))/\log r\to0$. Constants are handled by the convention above. The qualifier on the next scan is essential: $f(z)=2e^z$ fails $|f(0)|\le e^{0^k}$ for any positive $k$, although its order is finite. The manuscript has the eventual qualifier, so this is not a counterexample to its actual statement.

\section*{H0DT.11 --- Jensen's formula and reciprocal zero powers}
\textbf{Derivation, ours.} Assume $f(0)\ne0$. In a disk whose boundary contains no zero, divide $f$ by the finite product of its zero factors. The logarithm of the modulus of the remaining zero-free function is harmonic. The mean of $\log|re^{it}-a|$ is $\log\max(r,|a|)$: for $r<|a|$ expand $\log(1-re^{it}/a)$, and for $r>|a|$ factor $re^{it}$ and expand the reciprocal ratio. Taking real parts and means gives
\[
 V_f(r):=\frac1{2\pi}\int_0^{2\pi}\log|f(re^{it})|\,dt
 =\log|f(0)|+\sum_{|a_j|<r}\log\frac r{|a_j|}.
\]
At a boundary zero the logarithmic singularity is integrable and the same formula follows by continuity (the added term is zero). Thus, with $n(r)=\#\{j:|a_j|\le r\}$,
\[
 n(r)\log2\le V_f(2r)-\log|f(0)|\le\log M_f(2r)-\log|f(0)|.
\]
If $\rho<\infty$, for every $\delta>0$ this is $O(r^{\rho+\delta})$. For any $q>\rho$ choose $\delta<q-\rho$ and use the nonnegative identity
\[
 \sum_j|a_j|^{-q}=q\int_0^\infty n(t)t^{-q-1}\,dt,
\]
which follows by summing $q\int_{|a_j|}^\infty t^{-q-1}dt$. The integral is finite, since $n(t)=0$ near zero. Zeros at the origin are excluded exactly as the prime on the manuscript sum requires. The exact literature match is Valiron (1925), \S13, Theorem XXVI, p.~22, or McMullen (2025), Corollary 3.14, p.~86.

\section*{H0DT.12 --- Hadamard--Borel: full factorization proof}
The statement is Known: Valiron (1925), \S14, Theorem XXIX, p.~24, including the following definition of genus; McMullen (2025), Theorem 3.16, p.~89, gives the same factorization with a different convention for the word genus. Here is a complete proof of the needed assertion, ours.

Let $\tau=\inf\{s>0:\sum_j|a_j|^{-s}<\infty\}\le\rho$. If there are infinitely many zeros, let $p$ be the least nonnegative integer with $\sum_j|a_j|^{-p-1}<\infty$. Then $p\le\tau\le p+1$. Indeed, if $p>\tau$, the power $p$ would already be summable, contradicting minimality; the upper bound is immediate. Let $P=\prod_jE_p(z/a_j)$. For $p<q\le p+1$ one has
\[
 \log^+|E_p(w)|\le C_{p,q}|w|^q \quad(w\in\C).
\]
For $|w|\le1/2$ use H0DT.5 and $q\le p+1$; on $|w|\ge1/2$, use $\log(1+|w|)+\sum_{j=1}^p|w|^j/j\le C|w|^q$. If $\tau<p+1$, take any $q\in(\tau,p+1]$; if $\tau=p+1$, take $q=p+1$, whose power sum is finite by definition of $p$. Therefore $\ord P\le\tau$. Jensen's estimate applied to $P$ gives the reverse inequality. Thus $\ord P=\tau$.

We next give the lower modulus estimate needed for the exponential quotient. For every $\eta>0$ there are arbitrarily large radii $r$ with
\[
 \min_{|z|=r}\log|P(z)|\ge-C_\eta r^{\tau+\eta}. \tag{HB1}
\]
It suffices to consider $0<\eta<1$. From the zero exponent, $n(t)\le C_\delta t^{\tau+\delta}$ for large $t$ and every $\delta>0$ (use $n(t)t^{-q}\le\sum |a_j|^{-q}$ with $q=\tau+\delta$). Choose $0<\delta<\eta/4$ and $K>\tau+\delta+1$. In $[R,2R]$ exclude intervals of radius $R^{-K}$ around each modulus $|a_j|\le4R$. Their total length is $O(R^{\tau+\delta-K})<R$, so some $r$ remains. For $|z|=r$ and $|a_j|\le2r$,
\[
 \log|1-z/a_j|\ge-(K+1)\log R-C.
\]
There is a positive minimum zero modulus. Hence this logarithmic part, summed over the near zeros, is bounded below by $-C R^{\tau+2\delta}$. The real parts of the compensating polynomials have absolute sum at most
\[
 \sum_{j=1}^{p}\frac{r^j}{j}\sum_{|a_\nu|\le2r}|a_\nu|^{-j}
 \le C R^{\tau+2\delta}.
\]
To justify the last inequality explicitly, summation by parts with $n(t)\le Ct^{\tau+\delta}$ gives, for each $j\le p\le\tau$,
$\sum_{|a_\nu|\le2r}|a_\nu|^{-j}\le C(1+r^{\tau+\delta-j})$; after multiplying by $r^j$ this is $O(r^{\tau+\delta})$.
For the far zeros $|a_j|>2r$, H0DT.5 bounds the absolute logarithm by $2r^{p+1}\sum_{|a_j|>2r}|a_j|^{-p-1}$. If $\tau<p+1$, choose $q\in(\tau,\min(\tau+\eta/2,p+1))$. Then
\[
 r^{p+1}\sum_{|a_j|>2r}|a_j|^{-p-1}
 \le C r^q\sum_j|a_j|^{-q}=O(r^{\tau+\eta/2}).
\]
If $\tau=p+1$, the finite endpoint sum gives $O(r^\tau)$ directly. These estimates prove (HB1). For a finite zero family, $P$ is a polynomial and (HB1), with any positive exponent, is immediate on sufficiently large circles.

By Weierstrass, $f=z^m e^Q P$ with entire $Q$. On the chosen radii, finite order of $f$ and (HB1) imply
$\Re Q(z)\le C_\varepsilon r^{\rho+\varepsilon}$ for every $\varepsilon>0$ (the term $m\log r$ is harmless). Write $Q(z)=\sum_{n\ge0}q_nz^n$, and let $A_r$ be that upper bound. The mean of $\Re Q(re^{it})$ is $\Re q_0$. For $n\ge1$, Fourier coefficients give
\[
 |q_n|r^n=2\left|\frac1{2\pi}\int_0^{2\pi}\Re Q(re^{it})e^{-int}\,dt\right|
 \le2(A_r-\Re q_0).
\]
For any integer $n>\rho$, choose $0<\varepsilon<n-\rho$ and let these radii tend to infinity. Then $q_n=0$. Consequently $Q$ is a polynomial of degree at most $\rho$, while $p\le\tau\le\rho$. This proves genus $\max(p,\deg Q)\le\rho$ and the full factorization.

For clarity the usual equality $\rho=\max(\tau,\deg Q)$ also follows: the upper bound follows from the product and polynomial growth estimates. Jensen gives $\rho\ge\tau$. If $d=\deg Q>\tau$, use (HB1) in a fixed direction where the leading coefficient of $Q$ has positive real part; choosing $\eta<d-\tau$ gives $\log|f(re^{it})|\ge c r^d$ along the selected radii. Hence $\rho\ge d$. No lower-bound argument assumes the factorization theorem being proved.

\section*{H0DT.13 --- Order less than one}
If $\rho<1$, H0DT.12 gives the nonnegative integer genus $G\le\rho$, so $G=0$. Its product is precisely H0DT.9, with $z^m$ included when needed. This includes constants and polynomials under the stated convention.

\section*{H0DT.14 --- Order at most genus plus one}
If the genus is $G<\infty$, write $f=z^m e^Q\prod_jE_p(z/a_j)$ with $p\le G$, $\deg Q\le G$. The elementary bound from H0DT.12 with $q=p+1$ gives $\log M_P(r)\le C r^{p+1}\sum|a_j|^{-p-1}$. Also $\log|e^{Q(z)}|\le C(1+r^G)$ and $\log r^m=O(\log r)$. Hence $\rho\le G+1$. This proves the manuscript's unproved upper bound.

\section*{H0DT.15 --- Interchanged names in the boxed inequality}
\textbf{Literal counterexample.} If $\omega$ denotes genus and $g$ order, as the adjacent sentence says, take $f(z)=\prod_{n\ge1}(1-z/n^2)$. H0DT.18 proves genus $0$ and order $1/2$, so the written $\omega-1\le g\le\omega$ would demand $1/2\le0$. All zeros are nonzero and the function is entire. \textbf{Correction, ours.} Interchange the two nouns in that sentence, leaving the box intact: $\omega=\rho$, $g=G$, and $\rho-1\le G\le\rho$. The preceding two inequalities, proved in H0DT.12 and H0DT.14, provide direct local evidence for a Slip rather than a change of theorem.

\section*{H0DT.16 --- Integer and noninteger alternatives}
The integer $G$ lies in $[\rho-1,\rho]$. If $\rho$ is not an integer, $G=\lfloor\rho\rfloor$. If $\rho=q\ge1$ is an integer, $G\in\{q-1,q\}$. At $\rho=0$ only $G=0$ is possible; genus $-1$ is not permitted. Both alternatives for each $q\ge1$ occur. The function $e^{z^q}$ has genus $q$ and order $q$. For the other alternative set
\[
 a_n=\bigl(n[\log(n+1)]^2\bigr)^{1/q},\qquad
 f_q(z)=\prod_n E_{q-1}(z/a_n).
\]
Then $\sum|a_n|^{-q}<\infty$ by the integral test, while $\sum|a_n|^{-s}=\infty$ for $s<q$ (logarithms grow more slowly than every positive power). Thus the zero exponent is $q$ and the least primary degree is $q-1$. The product has order $q$ by the canonical-product part of H0DT.12 and genus $q-1$. These constructions prove sharpness, not merely compatibility of two integer bounds.

\section*{H0DT.17 --- The exponential example}
For an integer $q\ge1$, $M_{e^{z^q}}(r)=e^{r^q}$, so the order is $q$. The zero family is empty and the residual polynomial is $z^q$, so the genus is $q$. For $q=0$ the function is the nonzero constant $e$, with both invariants zero. The standard convention that $z^\omega$ is an entire monomial forces $\omega\in\N\cup\{0\}$; it is not a statement for arbitrary real powers. If that convention is omitted, $e^{z^{1/2}}$ does not define a single-valued entire function and cannot satisfy the premise. The parameter restriction in our precise statement makes this explicit.

\section*{H0DT.18 --- The product with roots $n^\alpha$}
\textbf{Refutation and exact correction.} Let $\alpha>1$ and $f(z)=\prod_{n\ge1}(1-z/n^\alpha)$. The reciprocal-root sum is $\sum n^{-\alpha}<\infty$, so it is entire of genus zero. Its zero exponent is $1/\alpha$, since $\sum n^{-\alpha s}$ converges exactly for $s>1/\alpha$. The canonical-product argument in H0DT.12 gives
\[
 \ord f=1/\alpha<1.
\]
This refutes the ink's order $1$, for every value of the stated parameter, and checks all hypotheses. An independent elementary upper estimate is in H0DT.31; Jensen supplies the matching lower bound. The most economical correction of this displayed example is $1\mapsto1/\alpha$. If its purpose was instead to illustrate order one with genus zero, the product with roots $n[\log(n+1)]^2$ in H0DT.16 does that. These are different corrections of different purposes; neither is silently attributed to his ink. The explicit reservation \emph{sauf erreur} is preserved.

\section*{H0DT.19 --- Sums, products, and algebras}
For entire $f,g$, $M_{fg}(r)\le M_f(r)M_g(r)$ and $M_{f+g}(r)\le M_f(r)+M_g(r)$. If $k>\max(\rho_f,\rho_g)$, both functions are eventually bounded by $e^{r^k}$, and the two new maxima are at most $e^{2r^k}$ and $2e^{r^k}$. Taking the double-logarithmic upper limit proves the asserted maximum bound; infinite orders cause no difficulty. Consequently the class of functions of order at most $\mu$ is a complex algebra for $\mu\ge0$, and the class of order strictly less than $\mu$ is one for $\mu>0$, after including the zero function. At $\mu=0$ the strict class, after adding the zero function, is the zero algebra; the usual nontrivial unital assertion uses $\mu>0$. Cancellation can lower order, so no equality is asserted.

\section*{H0DT.20 --- The coefficient criterion}
Assume $\alpha>0$, $f(z)=\sum_{n\ge0}c_nz^n$, and $n^\alpha|c_n|^{1/n}\to0$. Put $\beta=1/\alpha$. We prove the stronger estimate $\log M_f(r)=o(r^\beta)$, which implies order at most $\beta$ and, by H0DT.12, genus at most $\beta$.
For every $d>0$ there is $C_d$ such that $|c_n|\le C_d(d/n^\alpha)^n$ for all $n\ge1$. Set $t=(dr)^{1/\alpha}$. Then $|c_n|r^n\le C_d(t/n)^{\alpha n}$. The function $x\log(t/x)$ on $x>0$ has maximum $t/e$. For $n\le2et+1$, each term is at most $C_de^{\alpha t/e}$; there are at most $2et+2$ such terms. For $n>2et+1$, $(t/n)^{\alpha n}\le(1/(2e))^{\alpha n}$, a convergent geometric tail bounded uniformly in $t$. Therefore
\[
 \log M_f(r)\le (\alpha/e)d^{1/\alpha}r^{1/\alpha}+O_d(\log(r+2)).
\]
Letting $d\downarrow0$ proves the stronger claim. \textbf{Literal counterexample without positive exponent.} For $\alpha=-1$ and $f=e^z$, Stirling gives $n^{-1}(1/n!)^{1/n}\to0$, whereas genus $1$ is not at most $-1$. At $\alpha=0$ the displayed conclusion $1/\alpha$ is undefined. The surrounding positive exponent conventions justify a Slip, ours, rather than a change of mathematical intended claim.

\section*{H0DT.21 --- The coefficient-majorant proposition}
Let nonzero entire functions $f=\sum a_nz^n$, $h=\sum b_nz^n$ satisfy $|a_n|\le|b_n|$ for all $n$, with $h$ of genus zero. Suppose the nonzero zeros $a$ of $f$ satisfy $\sum |a|^{-1}<\infty$. Then $f$ has genus zero.
\textbf{Proof, ours.} Write $h=Cz^k\prod_j(1-z/b_j)$. Expanding absolute coefficients and passing to increasing finite products shows
\[
 \sum_n|b_n|r^n\le |C|r^k\prod_j(1+r/|b_j|)=H(r).
\]
For any $\varepsilon>0$ choose a finite head leaving reciprocal-root tail mass below $\varepsilon/2$. Its polynomial growth is bounded by a constant times $e^{\varepsilon r/2}$, and the tail product by $e^{\varepsilon r/2}$. Hence $H(r)\le C_\varepsilon e^{\varepsilon r}$. In particular $f$ has order at most one and is of exponential type zero. Hadamard--Borel with the minimal primary degree zero gives
\[
 f(z)=e^{a+bz}g(z),\qquad g(z)=z^m\prod_j(1-z/a_j).
\]
The same head-tail bound makes $g$ of type zero. Suppose $b\ne0$. Replacing the variable by $\overline b z/|b|$ reduces to $b>0$. Choose $\theta\in(\pi/3,\pi/2)$ and set $c=\cos\theta>0$. In $|\arg z|\le\theta$, the type-zero bound on $f$ with exponent $bc/2$ gives $|g(z)|\le C e^{-bc|z|/2}$. Set $d=bc/2$ and $\omega=e^{2\pi i/3}$. Among $z,\omega z,\omega^2z$, one lies in this sector. Each of the other two satisfies the global bound $|g(w)|\le C' e^{d|w|/3}$. Thus the entire function $G(z)=g(z)g(\omega z)g(\omega^2z)$ obeys $|G(z)|\le C''e^{-d|z|/3}$. It is bounded, hence constant, and its limit at infinity is zero. Therefore $G=0$, impossible since each factor is a nonzero entire function. So $b=0$, proving genus zero. This also proves the stronger equivalent criterion: a nonzero entire function of exponential type zero with summable reciprocal nonzero zeros has genus zero.

\section*{H0DT.22 --- The positive product majorant and the sum sign}
The construction required in H0DT.21 is $H(r)=|C|r^k\prod_j(1+|b_j|^{-1}r)$. Its coefficients are nonnegative and dominate the absolute coefficients of $h$ by the triangle inequality, first for finite products and then by uniform convergence on compact disks. Summability of the reciprocal roots proves that this is a genus-zero entire function. If a bare product with nonzero constant term is desired even when $k>0$, replace $r^k$ by $(1+r)^k$: this adds nonnegative coefficients and still gives a genus-zero product. Thus that harmless normalization does not remove any hypothesis. The literal $C\sum_j(1+\alpha_jz)$ written in place of a product is divergent already at $z=0$ for infinitely many factors and $C\ne0$, even when $\alpha_j=2^{-j}>0$ is summable. This checks every displayed condition and refutes that formula. The preceding genus-zero representation and the subsequent use of a product show the intended product; the Slip classification is ours.

\section*{H0DT.23 --- The rotation}
For $b\ne0$, $b(\overline b/|b|)=|b|$, and the multiplier has modulus one. Rotating the variable therefore makes the linear exponential coefficient positive and preserves absolute Taylor coefficients, zero moduli, order, and genus. The written $b/|b|$ instead gives $b^2/|b|$; for $b=i$ this is $-1$, not a positive real number. This is the complete counterexample to the literal substitution and proof of the conjugated correction.

\section*{H0DT.24 --- Genus zero implies exponential type zero}
Write $h=Cz^m\prod_j(1-z/a_j)$ with $\sum|a_j|^{-1}<\infty$. Given $\varepsilon>0$, choose $J$ with $\sum_{j>J}|a_j|^{-1}<\varepsilon/2$. Then
\[
 |h(z)|\le |C||z|^m\prod_{j\le J}(1+|z|/|a_j|)\exp(\varepsilon|z|/2).
\]
The finite polynomial in $|z|$ is at most $C_\varepsilon e^{\varepsilon|z|/2}$, since every polynomial times a decaying exponential is bounded on $[0,\infty)$. This proves the assertion for every $z$, with a constant depending on $\varepsilon$. This is stronger than order at most one and is not interchangeable with merely finite exponential type.

\section*{H0DT.25 --- The sector endpoint}
The literal range $\pi/3<\theta<\pi$ does not ensure $\Re z\ge c|z|$ with $c>0$ on $|\arg z|\le\theta$. Take $\theta=2\pi/3$ and $z=e^{2\pi i/3}$; then $\Re z=-1/2$. The corrected range $\pi/3<\theta<\pi/2$ gives $\Re z\ge|z|\cos\theta>0$. Its lower endpoint is needed because three sectors rotated by $2\pi/3$ must cover the circle. With $a\in\C$ and $b>0$ real after the preceding rotation, the exact lower estimate for $f=e^{a+bz}g$ is
\[
 |f(z)|=|e^a|e^{b\Re z}|g(z)|\ge |e^a|e^{b\cos\theta\,|z|}|g(z)|.
\]
The page also suppresses the constant and the factor $b$ in the rate. As a local algebraic inequality its coefficient-one form need not hold: take $a=-2$, $b=1$, $g=1$, and $z=1$, giving $|f(1)|=e^{-1}$, below $e^c|g(1)|$ for every $c>0$. This tests the intermediate exponential estimate, not the full proposition (whose hypotheses force $b=0$). Retaining $|e^a|$ and choosing the rate from $b\cos\theta$ gives exactly the constants used in H0DT.21. Thus both endpoints in the corrected choice have a stated role; the missing denominator is classified as a Slip on the local evidence of the proof, without changing the ink.

\section*{H0DT.26 --- Coefficient domination is radial}
For a concrete counterexample inside the relevant sector, let $h(z)=1+z$, $f(z)=1+e^{-i\pi/3}z$, and $z=e^{i\pi/3}$. Both functions are genus-zero polynomials with identical coefficient moduli; their reciprocal-zero sums are finite. Yet $|f(z)|=2>\sqrt3=|h(z)|$. Thus $|f(z)|\le|h(z)|$ does not follow. The strongest general elementary consequence is
\[
 |f(z)|\le\sum_n|a_n||z|^n\le\sum_n|b_n||z|^n,
\]
which equals $h(|z|)$ after the positive-coefficient replacement, and is bounded by the product majorant in H0DT.22 in general. That is exactly what the sector proof requires.

\section*{H0DT.27 --- The three-rotation argument}
The estimates and Liouville contradiction are proved in H0DT.21, independently of this manuscript step. The product identity used on the page is also exact: $\prod_{j=0}^2(1-\omega^jz/a)=1-z^3/a^3$ and $\omega^{3m}=1$. Absolute local convergence permits multiplication and regrouping, so
\[
 g(z)g(\omega z)g(\omega^2z)=z^{3m}\prod_i(1-z^3/a_i^3).
\]
The page's live word \emph{constante} is not replaced by \emph{zero} as a reading. Constancy follows by Liouville; zero follows separately from the displayed decaying upper bound. A nonzero constant $g$ cannot satisfy the sector decay either. This closes the argument also when there are no nonzero zeros, a case that a roots-only contradiction would miss.

\section*{H0DT.28 --- Equality of the two sequence exponents}
Let $a_n>0$ be nonincreasing with $a_n\to0$, and define
\[
 \rho=\inf\{s>0:a_n=O(n^{-1/s})\},\qquad
 \sigma=\inf\{s>0:\sum_n a_n^s<\infty\}.
\]
If $a_n\le C n^{-1/s}$, then for every $t>s$, $\sum a_n^t\le C^t\sum n^{-t/s}<\infty$. Thus $\sigma\le s$, and taking infima gives $\sigma\le\rho$. Conversely $\sum a_n^s<\infty$ implies $n a_n^s\le\sum_{j=1}^n a_j^s\le C$, so $a_n=O(n^{-1/s})$ and $\rho\le s$. Taking infima gives $\rho\le\sigma$. These implications also handle the values zero and infinity without assuming convergence at the boundary exponent. The stronger little-o result is proved separately below.

\section*{H0DT.29 --- The logarithmic exponent: limsup is wrong}
\textbf{Exact correction.} With the preceding hypotheses and $L_n=\log(1/a_n)/\log n$ for $n\ge2$,
\[
 \frac1\rho=\liminf_{n\to\infty}L_n,
 \qquad
 \rho=\limsup_{n\to\infty}\frac{\log n}{\log(1/a_n)},
\]
with reciprocal conventions $1/0=\infty$, $1/\infty=0$. To prove the first identity, $a_n=O(n^{-1/s})$ gives $\liminf L_n\ge1/s$. Conversely if $1/s<\liminf L_n$, then eventually $a_n\le n^{-1/s}$, so $s$ is admissible. These two implications prove equality, including endpoint values by taking strict approximations. Reciprocation of eventually positive numbers gives the second identity.

\textbf{Counterexample checking monotonicity.} Let $N_k=2^{2^k}$ and, for $N_k\le n<N_{k+1}=N_k^2$, set $a_n=N_k^{-2}$; fill the finite initial part with ones. The sequence is positive, nonincreasing, and tends to zero. In the $k$th block,
$L_{N_k}=2$, whereas $L_{N_{k+1}-1}\to1$, and all intermediate values lie between these endpoints. Hence $\liminf L_n=1$ and $\limsup L_n=2$. Moreover the block sum of $a_n^s$ is $(N_k^2-N_k)N_k^{-2s}$: these sums converge over $k$ if $s>1$ and diverge if $s\le1$. Thus $\rho=\sigma=1$, so the written formula $1/\rho=\limsup L_n$ is false. Changing an upper to a lower limit is our mathematical correction; the ink remains an upper limit. We classify this isolated reversal as a Slip, with lower confidence about authorial intention than about the counterexample.

\section*{H0DT.30 --- The sharp monotone little-o lemma}
If $a_n\ge0$ is nonincreasing and $\sum a_n^s<\infty$ for $s>0$, then for $n\ge2$,
\[
 (n-\lfloor n/2\rfloor)a_n^s
 \le\sum_{j=\lfloor n/2\rfloor+1}^n a_j^s\longrightarrow0.
\]
It follows that $n^{1/s}a_n\to0$. This is a complete independent proof of the strengthened live statement. The page's block argument also works: after rescaling a hypothetical positive limsup, choose increasing $n_i$ with $n_i a_{n_i}\ge1$; then $\varepsilon_i=(n_i-n_{i-1})/n_i$ has finite sum because the disjoint blocks of the convergent series dominate it. Starting with $i=2$, all $\varepsilon_i<1$, and $\prod_{i\ge2}(1-\varepsilon_i)>0$: eventually $\varepsilon_i\le1/2$ and $-\log(1-\varepsilon_i)\le2\varepsilon_i$. But $n_i=n_1/\prod_{j=2}^i(1-\varepsilon_j)$ would then be bounded, a contradiction. Starting at $i=2$ avoids the irrelevant first factor zero if an index $n_0=0$ is introduced. For $s\ne1$ apply this argument to $a_n^s$.

\section*{H0DT.31 --- The order of the positive product}
Let $a_n\ge0$, $\sum a_n<\infty$, and $F(z)=\prod_n(1+a_nz)$. Omit zero terms and arrange the rest in nonincreasing order; finite families give polynomials and exponent zero. Let $\sigma$ be their power-summability exponent. The zeros are $-1/a_n$, so Jensen in H0DT.11 gives $\sigma\le\ord F$. The reverse inequality follows from the canonical-product estimate in H0DT.12, or by the manuscript's comparison, which we now justify fully.

For $a_n=n^{-\beta}$, $\beta>1$, choose $\nu=\lfloor r^{1/\beta}\rfloor+1$. Since all coefficients are nonnegative, $M_F(r)=F(r)$, and
\[
 \log F(r)\le\nu\log(1+r)+r\sum_{n>\nu}n^{-\beta}
 \le\nu\log(1+r)+\frac{r\nu^{1-\beta}}{\beta-1}
 =O(r^{1/\beta}\log(r+2)).
\]
The tail comparison is $\sum_{n>\nu}n^{-\beta}\le\int_\nu^\infty t^{-\beta}dt$. This proves order at most $1/\beta$, while the zero exponent proves equality. For general $\sigma<1$, choose $s$ with $\sigma<s<1$. Since $\sum a_n^s<\infty$, monotonicity gives $a_n\le Cn^{-1/s}$, so comparison with this model at $Cr$ proves $\ord F\le s$. Let $s\downarrow\sigma$. If $\sigma=1$, H0DT.24 gives order at most one. The summability assumption excludes $\sigma>1$. Thus $\ord F=\sigma$ in every case, including zero and the endpoint one.

\section*{H0DT.32 --- Maximum term and order}
For an arbitrary entire $f(z)=\sum a_nz^n$ put $B(r)=\sup_n|a_n|r^n$. Cauchy gives $B(r)\le M_f(r)$. Conversely,
\[
 M_f(r)\le\sum_n|a_n|r^n\le B(2r)\sum_{n\ge0}2^{-n}=2B(2r).
\]
The logarithmic upper limits are unchanged by the constant factors and by replacing $r$ with $2r$. Therefore $\ord f=\limsup\log\log B(r)/\log r$, with the constant-function convention stated above. On scan 114 the coefficients are nonnegative, so his formula without modulus is valid in its local context. For complex coefficients the modulus is part of our extension, not an alteration of his preceding positive case.

\section*{H0DT.33 --- Replacing reciprocal roots by their moduli}
\textbf{Corrected statement and proof.} If $\sum|\lambda_j|<\infty$, then $f(z)=\prod_j(1+\lambda_jz)$ and $F(z)=\prod_j(1+|\lambda_j|z)$ are genus-zero entire functions with the same zero exponent $\sigma$. Jensen gives $\sigma\le\ord f$, the pointwise radial product inequality gives $M_f(r)\le F(r)$, and H0DT.31 gives $\ord F=\sigma$. Thus their orders coincide. This is the proof the page gives after its written word ``genre 1''.

\textbf{Why the literal genus-one generalization fails.} Enumerate $\lambda_{2n-1}=1/n$, $\lambda_{2n}=-1/n$. The partial bare products converge locally uniformly: paired products are $\prod_n(1-z^2/n^2)$, and an unpaired last factor tends uniformly to one. This entire function has order one and genus one, since its nonzero reciprocal-zero absolute sum is $2\sum1/n=\infty$ but the squares are summable. The positive replacement $\prod_n(1+|\lambda_n|r)$ diverges for every $r>0$; indeed $\log(1+r/n)\ge r/(2n)$ eventually. Thus it is not an entire function at all. The local product exists by the paired logarithmic series and has exactly the stated zeros. For an explicit order upper bound, put $N=\lceil r\rceil$. For $n\le N$, $1+r^2/n^2\le2N^2/n^2$, so the head logarithm is at most $N\log2+2N\log N-2\log N!=O(N)$ by integral comparison for $\log N!$. The tail is at most $r^2\sum_{n>N}n^{-2}\le r^2/N=O(r)$. Thus the order is at most one, while Jensen and the zero exponent give the lower bound one. The adjacent positive-product argument and the unregularized factors support genus zero as a Slip correction, ours.

\section*{H0DT.34 --- Order of an unordered scalar family}
For a countable indexed family $(z_i)$ define $\sigma=\inf\{s>0:\sum_i|z_i|^s<\infty\}$. A nonnegative sum is the supremum of its finite subsums, so it is unchanged by bijective reindexing. Therefore the exponent is well defined for an unordered family, including multiplicities; replacing a family by an ordinary set and deleting repeated values need not preserve it. For example the same value one repeated infinitely often has exponent infinity, whereas the one-point set has exponent zero. A finite family, or the all-zero family, has exponent zero. For an infinite family with any summable positive power, only finitely many terms exceed any fixed positive threshold and a nonincreasing rearrangement of its nonzero moduli exists. H0DT.28 then identifies this exponent with the manuscript's decreasing-sequence exponent. If no power is summable, the extended value infinity is well defined even when such a null rearrangement does not exist.

\section*{H0DT.35 --- Rapid decrease and exponent zero}
For a nonincreasing null sequence of nonnegative terms, exponent zero is equivalent to rapid decrease: $n^k a_n\to0$ for every $k>0$. If all positive powers are summable, apply H0DT.30 with $s=1/k$ (or a smaller positive exponent). Conversely rapid decrease gives $a_n^s=O(n^{-2})$ by taking $k=2/s$, hence every positive power is summable.

For an unordered family this means rapid decrease of the decreasing rearrangement. Without that convention the converse fails. Define $a_{2^{2^j}}=2^{-j}$ for $j\ge1$, and $a_n=e^{-n}$ at every other index. Every positive power is summable, because $\sum_j2^{-js}<\infty$ and $\sum_ne^{-ns}<\infty$. Yet $n a_n=2^{2^j-j}\to\infty$ along the special indices. Thus this enumeration is not rapidly decreasing. The preceding monotone framework and the immediate passage to unordered families support the rearrangement convention; its omission is a Slip, ours, not a theorem about every enumeration.

\section*{H0DT.36 --- The programme and its precise versions}
The page proposes two summability-based notions for sequences in Banach or Fr\'echet spaces, but does not define vector powers or a complete exponent scale. There is therefore no uniquely specified further conjecture to prove. The substantive claim it makes---that strong and absolute summability lead to distinct extensions---is validated by the two constructions in H0DT.37--38 and their explicit separating example. The programme itself is retained as \emph{programme} in the ledger; the precise versions are separate decided rows. The definitions and properties in the next sections are ours. They are not claimed to be a uniquely recoverable formula in the manuscript.

\section*{H0DT.37 --- Absolute summability order: a supported version}
Let $E$ be a Hausdorff Fr\'echet space and $(x_n)$ a sequence. Define
\[
 \sigma_a(x)=\inf\{p>0:\text{for every continuous seminorm }q,\ \sum_n q(x_n)^p<\infty\}.
\]
This is an extended nonnegative number, and its defining set is upward closed: for fixed $q$, summability at $p$ makes $q(x_n)\le1$ eventually, so powers $s>p$ are smaller on the tail. A continuous seminorm for an equivalent locally convex presentation is bounded by a constant times a maximum of finitely many generating seminorms; $\max_j b_j^p\le\sum_j b_j^p$ proves independence of presentation. Reindexing preserves every nonnegative sum. If $T:E\to F$ is continuous linear, composing a seminorm on $F$ with $T$ gives $\sigma_a(Tx)\le\sigma_a(x)$. For a Banach space it is enough to test its norm, since all continuous seminorms are norm-bounded; equivalent norms give the same exponent. For $E=\C$ this is H0DT.34. These facts establish well-definedness and the claimed extension, rather than merely naming it.

\section*{H0DT.38 --- Strong summability and a complete supported exponent scale}
For a continuous seminorm $q$ on $E$ put $B_q^\circ=\{\phi\in E':|\phi(x)|\le q(x)\text{ for all }x\}$. For $p>0$ define
\[
 W_{p,q}(x;N)=\sup_{\phi\in B_q^\circ}\left(\sum_{n>N}|\phi(x_n)|^p\right)^{1/p}.
\]
Our strong version is $\sigma_u(x)=\inf\{p>0:W_{p,q}(x;N)\to0\text{ for every }q\}$. It is an explicit completion of the programme, not a phrase read into the page. It agrees with scalar summability order, is invariant under reindexing and equivalent locally convex presentations, and cannot exceed $\sigma_a$. To prove these claims, use respectively the scalar dual unit disk, the supremum over all finite subsets outside a fixed finite set, domination of seminorms and their quotient dual norms by Hahn--Banach, and
$W_{p,q}(x;N)^p\le\sum_{n>N}q(x_n)^p$. For $s>p$, the scalar inequality $\|a\|_s\le\|a\|_p$ gives $W_{s,q}\le W_{p,q}$, so the defining set is upward closed. Continuous linear maps are contractive after pulling back seminorms. For detail on seminorm invariance, if $q\le C\max_{j\le d}q_j$, embed the quotient by $\bigcap\ker q_j$ into the finite product of the normed quotients with the maximum norm. Hahn--Banach extends any $q$-dominated functional to that product. It decomposes as $\sum_j\phi_j$, where each $\phi_j$ is bounded by $Cq_j$. For $p\ge1$, Minkowski bounds its tail by $C\sum_jW_{p,q_j}$; for $0<p<1$, the inequality $|\sum_jc_j|^p\le\sum_j|c_j|^p$ gives the analogous finite sum of $p$th powers. Each tends to zero. This proves the claimed independence, including the non-Banach exponent range. These are all the invariance properties asserted here.

At $p=1$ this is exactly unconditional (strong) summability in a complete locally convex space. Indeed Hahn--Banach and choice of complex phases give
\[
 W_{1,q}(x;N)=\sup_{F\subset\{n>N\}\text{ finite},\ |\eta_n|\le1}
 q\left(\sum_{n\in F}\eta_nx_n\right).
\]
If this tends to zero, the net of finite subsums is Cauchy for every seminorm, hence converges by completeness. Conversely a Cauchy net of finite subsums controls all subset sums outside a fixed finite head. Real coefficients in $[0,1]$ are integrals of indicators of nested subsets; splitting complex coefficients into four such parts bounds the phase-sum seminorm by four times the subset-sum bound. Thus the displayed tail tends to zero. Absolute summability implies this condition directly. This is the historical distinction used here; it is not convergence merely in the original enumeration.

\textbf{Strict separation, fully computed.} In $E=\ell^2$ let $x_n=e_n/n$. Then $\sigma_a=1$. For $0<p<2$, H\"older's inequality, with its equality case on finite subsets, gives
\[
 W_{p,\|\cdot\|_2}(x;N)=
 \left(\sum_{n>N}n^{-2p/(2-p)}\right)^{(2-p)/(2p)},
\]
where an infinite right side means an unbounded supremum. To check equality, for a finite subset choose $|y_n|^2$ proportional to $n^{-2p/(2-p)}$ in the dual unit ball, then increase the finite subset. This tail tends to zero exactly when $p>2/3$. For $p\ge2$, the scalar $\ell^p$ norm is bounded by the $\ell^2$ norm, giving $W_p(x;N)\le1/(N+1)$. Therefore $\sigma_u=2/3<1=\sigma_a$. In particular $\sum e_n/n$ converges unconditionally, since the norm squared of any tail subsum is at most $\sum_{n>N}n^{-2}$, but $\sum\|e_n/n\|=\infty$. No general equivalence of the two notions is asserted.

\section*{H0DT.39 --- Order of a bounded set by a closed balanced convex hull}
For either chosen sequence exponent $\sigma_*$ above define, explicitly as our completion,
\[
 o_*(B)=\inf\{s>0:\exists(x_n)\text{ with }\sigma_*(x)<s
 \text{ and }B\subset\overline{\aco}\{x_n:n\ge1\}\}.
\]
Using $\sigma_*(x)\le s$ instead of $<s$ gives the same infimum: any witness for $\le s$ is a witness for $<s+\varepsilon$. For a nonzero space the literal requirement $\sigma_*(x)=s$ also gives the same infimum. Given a witness of order below $s$, choose $v\ne0$ and interleave it with $n^{-1/s}v$. The new sequence has order exactly $s$ (both versions reduce to scalar order on this line), and its hull still contains the original one. Conversely an exact-$s$ witness is a $\le s$ witness. For the zero space, the exact-positive-order reading has no witness and its literal infimum is $+\infty$; our at-most-order convention assigns zero. Both definitions are well defined. This degenerate distinction is stated, not suppressed; no claim on the page chooses between them. The construction is a programme completion, and the order-bound convention is ours. Here $\aco$ is the absolutely convex hull: finite sums $\sum c_nx_n$ with $\sum|c_n|\le1$. The construction is well defined in $[0,\infty]$, is monotone under inclusion, and satisfies $o_*(TB)\le o_*(B)$ for continuous linear $T$, by applying $T$ to a witness and using continuity on its closure. These assertions follow directly from the proved invariance of $\sigma_*$ and the hull definition.

Every bounded subset of a finite-dimensional space has order zero: choose a basis and a bound $C$ on all coordinate moduli; the absolutely convex hull of the finitely many vectors $dC e_j$ contains the set. Append zeros to form an order-zero sequence. Conversely any finite-order witness is a null sequence. For the absolute version this follows from positive-power summability in every seminorm; for the strong version Hahn--Banach gives $q(x_n)\le W_{p,q}(x;n-1)\to0$. The closed absolutely convex hull of a null sequence in a Fr\'echet space is compact: for any convex balanced neighbourhood $U$, all tail vectors lie in $U$, their absolutely convex hull lies in $U$, and the finite head has compact absolutely convex hull. This proves total boundedness; closure and completeness give compactness. Thus a bounded set with noncompact closure, such as the unit ball of infinite-dimensional $\ell^2$, has order infinity under either construction. This proves that the extended-value convention is necessary. No claim that every compact set has finite order, or that the two bounded-set orders coincide, is imported into the manuscript.

\section*{H0DT.40 --- Finite open-interval counting identity}
Let $0<r_1\le r_2\le\cdots\to\infty$ (finite families are allowed), $n(t)=\#\{\nu:r_\nu\le t\}$, $\alpha>0$, and $0<a<b<\infty$. For each atom with $a<r_\nu<b$, integrate $\alpha\int_{r_\nu}^b t^{-\alpha-1}dt=r_\nu^{-\alpha}-b^{-\alpha}$. Atoms at or below $a$ contribute $a^{-\alpha}-b^{-\alpha}$ to the integral, and atoms at $b$ contribute zero. Summing, or integrating the step function, gives exactly
\[
 \sum_{a<r_\nu<b}r_\nu^{-\alpha}
 =b^{-\alpha}n(b-0)-a^{-\alpha}n(a+0)
 +\alpha\int_a^b n(t)t^{-\alpha-1}\,dt.
\]
At $a=0$, interpret the boundary term as its limit zero; $n$ vanishes in an initial interval. If the written $a\le b$ includes equality at an atom, take $r_\nu=\nu$, $a=b=1$, and any $\alpha>0$. The sum and integral are zero, but the boundary difference is $n(1-)-n(1+)=-1$. Thus $a<b$ is needed for the displayed endpoint convention; the strict interval $0<a<b$ on scan 118 supports this Slip correction. For $a=b$ a separately defined empty-interval identity is simply $0=0$, without unequal one-sided counts.

\section*{H0DT.41 --- Power sum and weighted counting integral}
For $\alpha>0$, Tonelli's theorem for nonnegative terms gives
\[
 \sum_\nu r_\nu^{-\alpha}
 =\sum_\nu\alpha\int_{r_\nu}^{\infty}t^{-\alpha-1}dt
 =\alpha\int_0^\infty n(t)t^{-\alpha-1}dt.
\]
The equality is valid in $[0,\infty]$, so convergence of either side is equivalent to convergence of the other. This proof does not assume the boundary limit it is subsequently used to prove. The case of finitely many radii follows by the same identity.

\section*{H0DT.42 --- Vanishing endpoint term}
Suppose the power sum in H0DT.41 is finite. For fixed $R>0$ and $r>R$,
\[
 0\le r^{-\alpha}n(r)
 \le r^{-\alpha}n(R)+\sum_{R<r_\nu\le r}r_\nu^{-\alpha}.
\]
First let $r\to\infty$ and then $R\to\infty$; the right side tends to zero. This proves $n(r)/r^\alpha\to0$ without monotone indexing beyond the counting definition. The manuscript's disjoint-block proof is valid after using a single name for its endpoints $b_i$ (which become $\lambda_i$ in the display); the quantitative split above provides a second route without that notational slip.

\section*{H0DT.43 --- The tail identity and its missing exponent}
Assume the power sum is finite. Let $b\to\infty$ through values avoiding atoms in H0DT.40 and use H0DT.42. For every $a\ge0$,
\[
 \sum_{r_\nu>a}r_\nu^{-\alpha}
 =-a^{-\alpha}n(a)+\alpha\int_a^\infty n(t)t^{-\alpha-1}\,dt,
\]
with the boundary convention at zero as before. The right side of the manuscript has precisely this power dependence; its left side omits $\alpha$. For $r_\nu=2^\nu$, $\nu\ge1$, $a=0$, and $\alpha=2$, the written left side is $\sum2^{-\nu}=1$, while the right side is $\sum4^{-\nu}=1/3$. Both converge and every hypothesis holds. This is an explicit counterexample to the literal formula and a complete proof of the corrected one.

\section*{H0DT.44 --- General decreasing weights}
Let $\lambda:(0,\infty)\to[0,\infty)$ be $C^1$, nonincreasing, and tend to zero; put $\mu=-\lambda'\ge0$. (Nonnegativity already follows from monotonicity and the zero limit.) Then $\lambda(t)=\int_t^\infty\mu(r)dr$. Tonelli gives the extended identity
\[
 \sum_\nu\lambda(r_\nu)=\int_0^\infty n(r)\mu(r)dr.
\]
If it is finite, split at $R$ exactly as in H0DT.42:
$\lambda(r)n(r)\le\lambda(r)n(R)+\sum_{R<r_\nu\le r}\lambda(r_\nu)$.
This proves $\lambda(r)n(r)\to0$. Finally
\[
 \sum_{r_\nu>a}\lambda(r_\nu)
 =\int_a^\infty[n(r)-n(a)]\mu(r)dr
 =-\lambda(a)n(a)+\int_a^\infty n(r)\mu(r)dr.
\]
There is no additional regularity assumption at zero because $n$ vanishes near zero. The same finite-endpoint identity holds with $\lambda(b)n(b-)-\lambda(a)n(a)+\int_a^bn\mu$. In fact local absolute continuity suffices for the equalities; the $C^1$ hypotheses are the page's convenient case.

\section*{H0DT.45 --- The circle count}
For $f(0)\ne0$ and a radius $r$ whose circle contains no zero, the residue of $f'/f$ at a zero of multiplicity $m$ is $m$. The residue theorem and $z=re^{i\phi}$ give
\[
 n(r)=\frac1{2\pi i}\int_{|z|=r}\frac{f'(z)}{f(z)}dz
 =\frac1{2\pi}\int_0^{2\pi}\frac{f'(re^{i\phi})}{f(re^{i\phi})}re^{i\phi}\,d\phi.
\]
At a zero radius the ordinary contour integral is undefined. For example $f(z)=1-z$, $r=1$, has a pole on the contour, although $n(1)=1$ is defined. Thus the omitted exceptional-radius convention is necessary and is a Slip only at the level of literal ``every $r$''. If $f$ has a zero at zero, this contour counts it as well; either include its multiplicity or divide by $z^m$. All later weighted nonzero-zero formulae use the latter convention.

\section*{H0DT.46 --- Why Fubini is legitimate on an annulus}
Fix $0<a<b<\infty$ and a bounded continuous weight $\mu$ there. There are finitely many zeros in a neighbourhood of the closed annulus. Locally at such a zero $z_0$, $f'/f=m/(z-z_0)+h(z)$ with holomorphic $h$. The function $1/|z-z_0|$ is area-integrable on a disk, since $\int_0^\delta t^{-1}t\,dt<\infty$. Polar measure $dr\,d\phi$ on this annulus is equivalent to area measure, with bounded Jacobian factors. Hence the manuscript's double integral of the logarithmic derivative times $r\mu(r)e^{i\phi}$ is absolutely integrable, and Fubini applies. The finitely many radii and rays meeting zeros cause no failure of this two-dimensional conclusion. Likewise $\log|z-z_0|$ is locally integrable, which justifies the subsequent real logarithmic means. An integration-by-parts formula involving $(r\mu)'$ additionally requires that weight derivative; it is explicitly imposed in H0DT.48.

\section*{H0DT.47 --- Principal logarithm versus the real-part annotation}
\textbf{Literal failure of the complex primitive step.} Take the zero-free entire function $f(z)=e^{iz}$, radial angle zero, $a=\pi/2$, $b=3\pi/2$, and $\mu(r)=1/r$, so $r\mu(r)=1$. Then
\[
 \int_a^b\frac{f'(r)}{f(r)}dr=i\pi,
\]
whereas the principal logarithms at the endpoints differ by $-i\pi$. At the intermediate crossing of the negative real axis the principal logarithm jumps by $-2\pi i$. Thus using the principal determination as a single differentiable radial primitive is false even with no zeros; all the finite-annulus and smooth-weight conditions hold.

\textbf{Repair already indicated on the page.} The oblique words instruct taking real parts. On each zero-free radial segment, $d\log|f(re^{i\phi})|/dr=\Re(e^{i\phi}f'/f)$. Imaginary branch jumps disappear under the real part. Exceptional rays have measure zero, and the logarithmic singularities are integrable by H0DT.46. Equivalently, one can bypass branch choices entirely with Jensen's locally absolutely continuous mean, as in H0DT.48. We classify the principal-Log intermediate formula as a Slip in a proof whose final real formula is correct, and retain the annotation rather than refuting a truncated version of the supplied passage.

\section*{H0DT.48 --- The circular mean and weighted integration by parts}
On every compact subinterval of $(0,\infty)$, Jensen's formula in H0DT.11 expresses $V(r)$ as a constant plus a finite sum of $\log^+(r/|a_j|)$, plus smooth contributions from smaller zeros. Therefore $V$ is locally absolutely continuous and $V'(r)=n(r)/r$ almost everywhere. If $\mu$ is $C^1$ on $[a,b]$, ordinary integration by parts yields
\[
 \int_a^b n(r)\mu(r)dr
 =\int_a^b r\mu(r)V'(r)dr
 =[r\mu(r)V(r)]_a^b-\int_a^b V(r)(r\mu(r))'dr.
\]
This proves the manuscript's identity at every pair of endpoints, including zero radii, without ever assigning a principal complex logarithm across a jump. Local absolute continuity of $r\mu$ suffices more generally.

\section*{H0DT.49 --- The finite weighted zero identity}
For a $C^2$ weight $\lambda$ on $[a,b]$, combine the finite identity in H0DT.44 with H0DT.48 for $\mu=-\lambda'$. The result is
\[
 \sum_{a<r_\nu<b}\lambda(r_\nu)
 =\lambda(b)n(b-)-\lambda(a)n(a)
 -[r\lambda'(r)V(r)]_a^b
 +\int_a^b V(r)\,d(r\lambda'(r)).
\]
The last integral is an ordinary integral against $(r\lambda')'dr$ in this smooth setting. The plus sign follows from $r\mu=-r\lambda'$ and agrees with the ink. Unlike the convergence tests, this finite identity does not require a nonnegative or decreasing weight; smoothness on the finite interval is enough. This is the strongest version proved here; the decreasing case is the one carried forward on the page.

\section*{H0DT.50 --- The lower endpoint zero}
Assume $|f(0)|=1$ and $f(0)\ne0$. Below its first positive zero modulus, Jensen gives $V(r)=\log|f(0)|=0$ and $n(r)=0$ exactly. Hence all lower-endpoint terms in H0DT.49 vanish for every sufficiently small $a$, irrespective of the rate at which $r^2\mu(r)$ tends to zero. Therefore
\[
 \sum_{r_\nu<b}\lambda(r_\nu)
 =\lambda(b)n(b-)-b\lambda'(b)V(b)
 +\int_0^b V(r)\,d(r\lambda'(r)).
\]
It is valid for weights with the local smoothness used above, even if they are singular at zero, because the integrands are identically zero initially. The additional $r^2\mu(r)\to0$ condition written on the page is sufficient but unnecessary. The exact vanishing is stronger than the elementary $V(r)=O(r)$ estimate one might otherwise use. The zero-free case gives $V=0$ everywhere and is included.

\section*{H0DT.51 --- What the maximum-modulus bound actually says}
Under the immediately preceding normalization $|f(0)|=1$, Jensen gives $0\le V(r)\le\log M_f(r)$. Since $M_f(r)\ge1$ and $\log x\le x$ for $x\ge1$, the ink's weaker bound $|V(r)|\le M_f(r)$ is true. No insertion of a missing $\log$ is needed. If read without that normalization, the absolute-value assertion is false: the constant entire function $f=e^{-2}$ has $V=-2$ and $M=e^{-2}$, hence $|V|>M$. The strongest general two-sided bound is
\[
 \log|f(0)|\le V(r)\le\log M_f(r),\qquad f(0)\ne0,
\]
proved in H0DT.11. Thus the status of the local normalized assertion is True as written; the scope warning and the unnormalized counterexample are kept beside it. This preserves both the parenthetical qualification and the actual symbol $M$. To make the announced weighted estimate explicit, for $\mu\in C^1[a,b]$ the identity in H0DT.48 and $0\le V\le L:=\log M$ give
\[
 \left|\int_a^b n(r)\mu(r)\,dr\right|
 \le |b\mu(b)|L(b)+|a\mu(a)|L(a)
       +\int_a^b L(r)|(r\mu(r))'|\,dr.
\]
Replacing $L$ by $M$ gives the weaker bound directly justified by the written inequality. When $\mu\ge0$, the negative lower boundary term can be discarded for a sharper upper bound. Thus the announcement itself, not only its preparatory comparison, is proved.

\section*{H0DT.52 --- Elementary symmetric coefficients and the canonical product}
Let $\lambda\in\ell^1(\C)$, put $\alpha_0=1$, and for $n\ge1$ define
$\alpha_n(\lambda)=\sum_{i_1<\cdots<i_n}\lambda_{i_1}\cdots\lambda_{i_n}$.
Expanding $(\sum_i|\lambda_i|)^n$ shows that the contribution of distinct ordered tuples is $n!\sum_{i_1<\cdots<i_n}|\lambda_{i_1}\cdots\lambda_{i_n}|$, so the defining series is absolutely convergent and bounded by $\|\lambda\|_1^n/n!$. The finite-product expansions converge uniformly on $|z|\le R$, since their absolute sums are bounded by $e^{R\|\lambda\|_1}$ and the tail factors converge uniformly. Hence
\[
 \prod_i(1+\lambda_i z)=\sum_{n\ge0}\alpha_n(\lambda)z^n.
\]
The nonzero zeros are $-1/\lambda_i$ and their reciprocal moduli are summable; the product has genus zero and value one at zero. Every finite-support case is included. The doubtful argument of $f$ in the ink is not used to infer an extra dependence: $z$ is the independent variable in this precise reconstruction.

\section*{H0DT.53 --- The missing modulus}
For complex $\lambda$, $\alpha_n(\lambda)\le\|\lambda\|_1^n/n!$ is not an ordered real inequality: take $\lambda=(i,0,\ldots)$ and $n=1$, when its left side is $i$. The corrected statement is
\[
 |\alpha_n(\lambda)|\le\alpha_n(|\lambda|)\le\|\lambda\|_1^n/n!.
\]
The first inequality is termwise triangle inequality; the second is the distinct-tuple expansion proved in H0DT.52. It is sharp as a supremum on the unit ball, as H0DT.65--66 show. Restricting to nonnegative coordinates would also repair the display, but would weaken the surrounding complex proposition unnecessarily; adding the modulus is the evident Slip correction, ours.

\section*{H0DT.54 --- Which genus-zero coefficients are obtained}
By H0DT.9, every nonzero genus-zero function with nonzero constant term is $C\prod_i(1+\lambda_i z)$ with $\lambda\in\ell^1$. Its coefficients are exactly $C\alpha_n(\lambda)$, so the manuscript's assertion holds in the standing nonvanishing-at-zero convention. Without it, $f(z)=z$ is genus zero but cannot be a nonzero constant times such a product, whose value at zero is always that constant. The full statement is $Cz^m\prod_i(1+\lambda_i z)$: its coefficients are zero below $m$ and $C\alpha_{n-m}(\lambda)$ at $n\ge m$. The statement is classified as a Slip when read without the convention; the evidence for the intended scope is the assumption on scan 107 and the exact genus-zero formula on scan 108. No root at zero is counted in $\lambda_i=-1/a_i$.

\section*{H0DT.55 --- Absolute convergence of $S$ at every point}
For $x_i=|\lambda_i|$, choose $m$ so that $\delta=\sum_{i>m}x_i<1$. For $t\ge0$,
\[
 \prod_i(1+tx_i)\le\prod_{i\le m}(1+tx_i)e^{\delta t}.
\]
Tonelli's theorem and $\int_0^\infty e^{-t}t^n dt=n!$ (by $n$ integrations by parts) give
\[
 \sum_{n\ge0}n!|\alpha_n(\lambda)|
 \le\sum_{n\ge0}n!\alpha_n(x)
 =\int_0^\infty e^{-t}\prod_i(1+tx_i)dt<\infty,
\]
because a fixed polynomial times $e^{-(1-\delta)t}$ is integrable. Thus $S(\lambda)=\sum n!\alpha_n(\lambda)$ is defined for every $\lambda\in\ell^1$, and dominated convergence also gives its complex integral representation with $x_i$ replaced by $\lambda_i$. This proof is independent of the faulty exponent in the manuscript's Cauchy estimate.

\section*{H0DT.56 --- The radius at the origin is exactly one}
Put $P_n(\lambda)=n!\alpha_n(\lambda)$. H0DT.53 gives $\sup_{\|\lambda\|_1\le1}|P_n(\lambda)|\le1$. For $m\ge n$ take $a_i^{(m)}=1/m$ for $i\le m$ and zero otherwise. Then
\[
 P_n(a^{(m)})=\frac{m(m-1)\cdots(m-n+1)}{m^n}\longrightarrow1
 \quad(m\to\infty).
\]
Therefore every homogeneous polynomial norm is exactly one. For any $r<1$, $\sum\sup_{\|\lambda\|\le r}|P_n(\lambda)|\le\sum r^n<\infty$. For $r\ge1$ the homogeneous terms do not even tend to zero uniformly, since their suprema are $r^n$. Thus the radius of uniform convergence on centred balls is exactly one. H0DT.65 gives the same radius for the associated multilinear norms. This is compatible with pointwise convergence everywhere and does not assert a singularity of the function on every point of the sphere.

\section*{H0DT.57 --- Global holomorphy, hence all finite-dimensional restrictions}
We prove a stronger statement than the proposition's restrictions: $S$ is Fr\'echet-holomorphic on all of $\ell^1$. Fix $\lambda^0$ and choose a finite head of length $m$ leaving tail norm below $1/4$. On a sufficiently small norm neighbourhood of $\lambda^0$, the tail norm is at most some $\delta<1$, and the head norm is bounded by $K$. Let $e_j$ denote the elementary symmetric sums of the absolute head coordinates. They are bounded uniformly, for example by $K^j/j!$. The convolution of head and tail coefficients gives
\[
 n!|\alpha_n(\lambda)|\le
 \sum_{j=0}^{\min(m,n)} e_j\frac{n!}{(n-j)!}\delta^{n-j}.
\]
The series over $n$ of this bound converges uniformly because it is bounded by a constant times a polynomial in $n$ times $\delta^{n-m}$. Hence the series of continuous homogeneous polynomials $P_n$ converges uniformly on that neighbourhood. Such a limit is holomorphic: on each smaller ball, scalar Cauchy integrals on complex lines bound the first and second derivatives of each tail by constants times its uniform supremum, so the derivatives converge in operator norm and the remainder is uniformly quadratic. This proves Fr\'echet differentiability directly. In particular its restriction to any finite-dimensional complex subspace is holomorphic everywhere. The margin on scan 122 asserts the global conclusion, and this argument supplies the missing local boundedness justification.

\section*{H0DT.58 --- Unbounded on the unit ball, but globally holomorphic}
For $a^{(m)}$ from H0DT.56 and every fixed $L$, positivity gives
\[
 S(a^{(m)})\ge\sum_{n=0}^L P_n(a^{(m)})\longrightarrow L+1.
\]
Since $L$ is arbitrary, $S(a^{(m)})\to\infty$. The norms are one. Even the open unit ball is unbounded for $S$: use $b^{(m)}=(1-1/m)a^{(m)}$ for $m\ge2$; for each fixed $n$, $P_n(b^{(m)})\to1$, so the same finite-partial-sum argument applies.

Grothendieck (1956), Introduction \S2, p.~324, equations (13)--(14), defines an \emph{analytique enti\`ere} map by a homogeneous expansion whose multilinear norms have $n$th roots tending to zero; the following paragraph gives absolute uniform convergence on every bounded ball. Our $P_n$ have multilinear norms one, by H0DT.65. Thus $S$ fails that historical definition, exactly as the manuscript says. In modern terminology it is an entire holomorphic function which is not of bounded type. This is a translation of terminology, not a counterexample to his assertion. Conversely, for globally holomorphic maps on a Banach space, boundedness on each ball implies decay of the polynomial norms by Cauchy at arbitrary radius; polarization changes those norms by at most $n^n/n!$, whose $n$th root stays bounded. Thus the historical condition corresponds to bounded type, not just local holomorphy.

\section*{H0DT.59 --- Factorial decay of genus-zero coefficients}
Let $f(z)=\sum c_nz^n$ be of genus zero, allowing a factor $z^m$. By H0DT.24, for every $\varepsilon>0$, $M_f(R)\le C_\varepsilon e^{\varepsilon R}$. Cauchy with $R=n/\varepsilon$ yields
\[
 |c_n|\le C_\varepsilon(e\varepsilon/n)^n.
\]
The elementary factorial estimate $\log n!=n\log n-n+O(\log n)$, obtained by comparing $\sum_{j=1}^n\log j$ with $\int_1^n\log t\,dt$, gives $(n!)^{1/n}\sim n/e$. Consequently
$\limsup(n!|c_n|)^{1/n}\le\varepsilon$. Letting $\varepsilon\downarrow0$ proves the limit zero. Equivalently $|c_n|^{1/n}=o(1/n)$. Choose any complex $n$th root of $n!c_n$ (zero if $c_n=0$); its modulus tends to zero, so $c_n=\varepsilon_n^n/n!$ with $\varepsilon_n\to0$. This proves the full coefficient assertion, including zero coefficients and general complex phases.

\section*{H0DT.60 --- Homogeneity and factorial equivalences}
Each defining monomial gives $\alpha_n(z\lambda)=z^n\alpha_n(\lambda)$; absolute convergence permits summing this equality. Since $z\lambda\in\ell^1$ for every $z$, H0DT.55 shows $\sum_n n!\alpha_n(\lambda)z^n$ converges absolutely for every complex $z$. The scalar root test gives $(n!|\alpha_n(\lambda)|)^{1/n}\to0$. Conversely that limit gives convergence for every $z$ by the same test. The relation $(n!)^{1/n}\sim n/e$, proved at the strength actually needed in H0DT.59, proves equivalence with $|\alpha_n|^{1/n}=o(1/n)$ and with the $\varepsilon_n$ representation. Thus the logical passage on scan 121 is sound independently of the subsequent erroneous coefficient estimate.

\section*{H0DT.61 --- Cauchy's exponent}
For $f(z)=\sum c_nz^n$, Cauchy's formula is
\[
 c_n=\frac1{2\pi i}\int_{|z|=R}\frac{f(z)}{z^{n+1}}dz,
 \qquad |c_n|\le M_f(R)/R^n.
\]
The contour length $2\pi R$ supplies the missing factor $R$. The literal $R^{n+1}$ bound is false: $f=1+z$, $n=1$, $R=2$ would give $1\le3/4$. If one insists on the proof's strictly positive infinite sequence, take $\lambda_1=1$, $\lambda_i=2^{-i-4}$ for $i\ge2$. It is positive, nonincreasing, summable, and has tail mass $1/32$. Then $\alpha_1=33/32$, while $M_f(2)\le3e^{1/16}<16/5$, so $M_f(2)/4<4/5<\alpha_1$. Here $e^{1/16}<1/(1-1/16)=16/15$ follows termwise from the exponential and geometric series. This verifies the counterexample even under all the reduction's strict hypotheses.

\section*{H0DT.62 --- Tail estimate and the exponential factor}
For nonnegative summable $\lambda_i$, define $\rho_N=\sum_{i>N}\lambda_i$. Since $1+t\le e^t$ for $t\ge0$,
\[
 \prod_{i>N}(1+R\lambda_i)\le e^{R\rho_N}.
\]
With the page's $N=n+1$ and $R=(n+1)/\rho_{n+1}$, this is $e^{n+1}$, not $e^n$. A fully positive counterexample to the latter tail bound is $n=1$, $\lambda_1=\lambda_2=\lambda_3=1$, $\lambda_{3+j}=2^{-j}/100$ for $j\ge1$. Then $\rho_2=101/100$, $R=200/101$, and the tail product is at least $1+R=301/101>e$. For instance $e<11/4<301/101$, with the first bound obtained from the exponential series and a geometric bound on its tail. The corrected exponential follows immediately from the displayed substitution. If $\rho_N=0$, the function is a polynomial and division by $\rho_N$ must be avoided; all sufficiently high coefficients vanish. That finite-tail convention is made explicit in our repaired proof.

\section*{H0DT.63 --- A consistent repair of the epsilon-product computation}
Assume first that $\lambda_i>0$ are nonincreasing and summable, and fix $N\ge1$. The correct Cauchy estimate and an $N$-factor head give, when $\rho_N>0$,
\[
 \alpha_N\le R^{-N}\prod_i(1+R\lambda_i)
 \le\prod_{j=1}^N(R^{-1}+\lambda_j)e^{R\rho_N}.
\]
Choose $R=N/\rho_N$. Since $j\le N$ and $\rho_j\ge\rho_N$,
\[
 \alpha_N\le e^N\prod_{j=1}^N(\rho_N/N+\lambda_j)
 \le\frac{e^N}{N!}\prod_{j=1}^N(\rho_j+j\lambda_j).
\]
Let $\delta_j=\rho_j+j\lambda_j$. Summability gives $\rho_j\to0$, and H0DT.30 gives $j\lambda_j\to0$, hence $\delta_j\to0$. Thus
\[
 (N!\alpha_N)^{1/N}\le e\left(\prod_{j=1}^N\delta_j\right)^{1/N}\longrightarrow0.
\]
If the sequence is finite, sufficiently high coefficients vanish; for complex sequences dominate by absolute values and rearrange their moduli. This proves the intended estimate with every index specified. It replaces the whole inconsistent $n+1$-factor chain as our reconstruction; the literal chain is refuted already by H0DT.61, and its exponential substitution separately by H0DT.62. No uncertain epsilon subscript is resolved by selecting the value this proof would prefer.

\section*{H0DT.64 --- Geometric means and the continuation}
If $\delta_j\ge0$ and $\delta_j\to0$, then $(\prod_{j=1}^N\delta_j)^{1/N}\to0$. If any factor is zero this is immediate. Otherwise, for any $\varepsilon>0$ choose $J$ with $\delta_j\le\varepsilon$ for $j>J$. The product root is at most $(\prod_{j\le J}\delta_j)^{1/N}\varepsilon^{1-J/N}$, whose limit superior is at most $\varepsilon$. Let $\varepsilon\downarrow0$. Applying this to H0DT.63 completes the corrected coefficient proof. A fixed positive factor $(\varepsilon_0/(n+1))^{1/n}$ tends to one; if the factor is zero it only strengthens the conclusion. The conclusion $(n!|\alpha_n|)^{1/n}\to0$ then implies absolute convergence by the root test. Scan 122 really continues the proof; the boundary after ``D'o\`u'' on scan 121 is not an authorial admission of an unresolved problem.

\section*{H0DT.65 --- The symmetric multilinear form and its norm}
For $n\ge1$ define
\[
 A_n(x^{(1)},\ldots,x^{(n)})=\frac1{n!}
 \sum_{i_1,\ldots,i_n\ {
m distinct}}
 x^{(1)}_{i_1}\cdots x^{(n)}_{i_n}.
\]
The sum converges absolutely, being bounded before the $1/n!$ factor by $\prod_{j=1}^n\|x^{(j)}\|_1$. It defines a continuous symmetric $n$-linear form; symmetry follows by permuting the tuple of distinct indices. On the diagonal every unordered index set is counted $n!$ times, so $A_n(\lambda,\ldots,\lambda)=\alpha_n(\lambda)$. Its norm is at most $1/n!$ and in fact equals $1/n!$: on the distinct unit vectors $e_1,\ldots,e_n$, exactly one ordered tuple contributes and the value is $1/n!$. Thus the multilinear norm of $n!A_n$ is one, agreeing with the polynomial norm calculation of H0DT.56. The scalar bound $|\alpha_n|\le\|\lambda\|_1^n$ recalled in the continuation is a weaker consequence of the sharper factorial bound, not a contradictory estimate.

\section*{H0DT.66 --- Equal coordinates and both radii}
For $a^{(m)}=(1/m,\ldots,1/m,0,\ldots)$ with $m$ nonzero entries, $\|a^{(m)}\|_1=1$ and
\[
 \alpha_n(a^{(m)})=\binom mn m^{-n}\quad(0\le n\le m),\qquad
 \alpha_n(a^{(m)})=0\quad(n>m).
\]
Multiplying by $n!$ gives the product of $1-j/m$, $j=0,\ldots,n-1$, tending to one for fixed $n$. H0DT.58 proves unboundedness by finite partial sums. H0DT.65 gives strict (multilinear-norm) radius $R_0=1$; H0DT.56 gives the polynomial-norm or large radius $R_1=1$. This verifies both inequalities and the equality on the page. It also checks the finite example in the model note: for four coordinates $1/8$,
\[
 S=\sum_{n=0}^4n!\binom4n8^{-n}
 =1+\frac12+\frac3{16}+\frac3{64}+\frac3{512}
 =\frac{891}{512}.
\]
An exact rational computation independently confirms this in the recorded script.

\section*{H0DT.67 --- Mellin integral for a positive product}
Let $\alpha_i\ge0$ and $0<p<1$. For real $r>0$ define $\log f(r)=\sum_i\log(1+\alpha_i r)$ as an extended nonnegative sum; an entire complex product exists if $\sum\alpha_i<\infty$. The extended real identity requires no such extra condition. Integration by parts for one positive $a$, with vanishing boundary terms at zero and infinity, gives
\[
 \int_0^\infty\log(1+ar)r^{-1-p}dr
 =\frac{a^p}{p}\int_0^\infty\frac{t^{-p}}{1+t}dt
 =\frac{\pi a^p}{p\sin\pi p}.
\]
Here the last integral can be verified without a table: integrate $z^{-p}/(1+z)$ around a keyhole contour cut along the positive axis, with argument in $(0,2\pi)$. The two sides contribute $(1-e^{-2\pi ip})J$; the only pole is $-1$ with residue $e^{-\pi ip}$. The small and large circles vanish because $0<p<1$. Hence $(1-e^{-2\pi ip})J=2\pi i e^{-\pi ip}$ and $J=\pi/\sin\pi p$. For $a=0$ both sides are zero. Tonelli now yields
\[
 \sum_i\alpha_i^p=\frac{p\sin\pi p}{\pi}
 \int_0^\infty\frac{\log f(r)}{r^{1+p}}dr
\]
with either side allowed to be infinity, exactly as his parenthesis says. If $A$ is Hilbert trace class, substitute its singular values: $M_{|A|}(r)=\det(I+r|A|)=\prod_i(1+rs_i(A))$. This gives the corollary on scan 130 for $S_p(A)=\sum s_i(A)^p$. The cancellation of the power-summability assumption on scan 129 is mathematically legitimate because of the extended-value interpretation; it must not be silently restored.

\section*{H0DT.68 --- Jensen's power-sum bound and the eigenvalue application}
For nonzero entire $f$ with $f(0)=1$, Jensen and Tonelli give, for every $p>0$,
\[
 \int_0^\infty V_f(r)r^{-1-p}dr
 =\sum_j\int_{|a_j|}^\infty\log(r/|a_j|)r^{-1-p}dr
 =\frac1{p^2}\sum_j|a_j|^{-p}.
\]
The final one-variable integral is $|a_j|^{-p}/p^2$ by integration by parts or substitution followed by $\int_1^\infty\log t\,t^{-1-p}dt=1/p^2$. Because $V_f\le\log M_f$, this proves
\[
 \sum_j|a_j|^{-p}\le p^2\int_0^\infty\log M_f(r)r^{-1-p}dr,
\]
including infinite values. For $f(z)=\det(I\pm zA)$ the zero sum is $s_p(A)=\sum|\lambda_j(A)|^p$. If $0<p<1$ and $A$ is trace class, H0DT.2 and H0DT.67 give
\[
 s_p(A)\le\frac{p\pi}{\sin\pi p}S_p(A).
\]
Both signs in the determinant give the same zero moduli and circle maxima. The coefficient tends to one as $p\downarrow0$ because $\sin t/t\to1$.

For completeness we justify the page's limiting improvement, including the operator inequality it cites to a later section. Exterior powers give, for compact operators $A,B$,
\[
 \prod_{j=1}^k s_j(AB)=\|\Lambda^k(AB)\|
 \le\|\Lambda^k A\|\,\|\Lambda^k B\|
 =\prod_{j=1}^k s_j(A)s_j(B).
\]
We use the elementary log-majorization implication: if decreasing positive finite sequences $a,b$ satisfy $\sum_{j\le k}\log a_j\le\sum_{j\le k}\log b_j$ for each $k$, then $\sum a_j^t\le\sum b_j^t$ for every $t>0$. Indeed
$\sum_j(\log a_j-c)_+=\max_k(\sum_{j\le k}\log a_j-kc)$ is at most the corresponding maximum for $b$. Integrate against $t^2e^{tc}dc$ on the real line and use
$e^{tx}=t^2\int_{-\infty}^x(x-c)e^{tc}dc$.
Zero entries follow by limits. Apply this argument to the first $N$ terms of the singular-value inequality and then let $N\to\infty$. Scalar H\"older with $1/r=1/p+1/q$ gives
\[
 S_r(AB)\le\sum_j(s_j(A)s_j(B))^r
 \le S_p(A)^{r/p}S_q(B)^{r/q},
\]
valid also when any sum is infinite and for all positive exponents in the finite case. Iteration gives $S_{p/n}(A^n)\le S_p(A)$. If $A\in S_p$ and $n>p$, this makes $A^n$ trace class, since an $\ell^{p/n}$ sequence with $p/n<1$ is summable. Algebraic spectral multiplicity gives $s_{p/n}(A^n)=s_p(A)$; eigenvalues that coalesce under the power have their multiplicities added. Apply the previous determinant bound to $A^n$ with exponent $p/n$:
\[
 s_p(A)\le\frac{(p/n)\pi}{\sin((p/n)\pi)}S_{p/n}(A^n)
 \le\frac{(p/n)\pi}{\sin((p/n)\pi)}S_p(A).
\]
Let $n\to\infty$ to obtain $s_p(A)\le S_p(A)$. If $S_p(A)=\infty$ the inequality is automatic. This supplies the exact missing gate of the application without presuming the Schatten inequality being derived elsewhere in the folder. The earlier integer $n$ is allowed to vary; a single fixed power would not remove the constant.

\section*{H0DT.69 --- The earlier determinant inequality: a reading correction}
The native-resolution crop of scan 71, Corollary 4, shows $\le$, not the supplied transcription's $=$. This is a palaeographic correction recorded in READINGS.tsv. H0DT.2 at $z=1$ proves the written inequality, with every trace-class hypothesis checked there. Equality in the alternative transcription would fail even in one dimension: $A=[-1]$ is trace class, $|\det(I+A)|=0$, and $\det(I+|A|)=2$. This counterexample diagnoses the importance of the sign; it is not charged to the manuscript. The stronger radial inequality of H0DT.2 is the exact continuation on scan 79.

\section*{H0DT.70 --- A maximum of complex determinants needs a modulus}
For $A=[1]$ and $r=1$, the values of $\det(I+zA)=1+z$ on $|z|=1$ include $1+i$ and $1-i$. No ordered-real supremum of these complex numbers is defined. Thus the omitted absolute-value bars are material. The corrected maximum exists by continuity on a compact circle. For positive trace-class $A$, every coefficient is nonnegative, so $|\sum\alpha_n z^n|\le\sum\alpha_n r^n$, with equality at $z=r$; hence $M_A(r)=\det(I+rA)$. Continuity in $r$ follows from uniform continuity on any compact disk, or from the maximum over the fixed compact unit circle of $|d_A(re^{it})|$. The local evidence for the Slip is the absolute-value definition at scan 108 and the positive-operator specialization immediately following this definition.

\section*{H0DT.71 --- Determinant submultiplicativity for a sum}
We prove the coefficient inequality
\[
 \alpha_n(|A+B|)\le\sum_{k=0}^n\alpha_k(|A|)\alpha_{n-k}(|B|).
\]
First let $A,B$ be finite rank. Expand $\Lambda^n(A+B)$ by multilinearity and group terms choosing $k$ copies of $A$. A singular-value decomposition writes each $\Lambda^k A$ as a sum of rank-one operators on decomposable unit exterior vectors, with total coefficients $\alpha_k(|A|)$; similarly for $\Lambda^{n-k}B$. Under the normalized symmetric exterior product, a pair $u\otimes u'^*$ and $v\otimes v'^*$ contributes
\[
 \binom nk^{-1}(u\wedge v)\otimes(u'\wedge v')^*.
\]
Its trace norm is at most $\binom nk^{-1}$, because the exterior norm of the wedge of the two orthonormal lists is the square root of their Gram determinant, at most one (successive orthogonal projection, or Gram--Schmidt). The binomial factor in the expansion cancels this factor. The trace-norm triangle inequality therefore gives the claimed coefficient convolution, since $\|\Lambda^n C\|_1=\alpha_n(|C|)$.

For trace-class operators use finite-rank truncations in trace norm. Multilinearity and the bound $\|C_1\wedge\cdots\wedge C_n\|_1\le\prod_j\|C_j\|_1/n!$, proved by rank-one expansions and the same Gram estimate, show continuity of $\Lambda^n C$ in trace norm. Hence the finite-rank convolution passes to the limit. Multiply by $r^n$ and sum nonnegative terms to obtain
\[
 \det(I+r|A+B|)\le\det(I+r|A|)\det(I+r|B|),\qquad r\ge0.
\]
Taking real logarithms and applying H0DT.67 gives $S_p(A+B)\le S_p(A)+S_p(B)$ for $0<p<1$. Tonelli permits infinite values. This proves the application without using the Schatten inequality as an input to the determinant bound.

\section*{H0DT.72 --- The second Mellin denominator}
The first integral on scan 131 has $r^{p+1}$; the next has a short minus-like stroke in $r^{p-1}$, followed by an appeal to formula (9). Retaining the latter reading, the identification needed to obtain the next finite coefficient is false. Take the trace-class scalar $A=[1]$, $p=1/2$. Then
\[
 \int_0^\infty\frac{\log(1+r)}{r^{p-1}}\,dr=\infty,
 \qquad
 \int_0^\infty\frac{\log(1+r)}{r^{p+1}}\,dr=2\pi.
\]
The first divergence follows already from $\log(1+r)r^{1/2}\ge r^{1/2}\log r$ for $r\ge1$; the second value is H0DT.67. Thus the infinity-valued inequality alone is true but cannot be the asserted application of (9). Restoring $p+1$ makes both the Jensen inequality and the Mellin substitution in H0DT.68 valid, giving the coefficient $p\pi/\sin(\pi p)$. The mathematical correction is certain; the confidence in the tiny original sign remains separately medium.

\section*{H0DT.73 --- Positive maximum comparison and Mellin change of variables}
Let $A\in S_p$, $B\in S_q$, $p,q>0$, $1/r=1/p+1/q$, and take an integer $m>\max(p,q)/2$. Put $H=A^*A$, $K=BB^*$, $L=(AB)(AB)^*$. The ideal inequality $s_j(XY)\le\|Y\|s_j(X)$ follows from the minimum error of rank-less-than-$j$ approximations: approximate $X$ and multiply the approximation by $Y$. Consequently, for $t>0$,
\[
 \sum_js_j(XY)^t\le\|Y\|^t\sum_js_j(X)^t
 \le\left(\sum_js_j(X)^t\right)\left(\sum_js_j(Y)^t\right).
\]
Apply this with $X=\Lambda^n A$, $Y=\Lambda^n B$, $t=2m$. The exterior spectral identities give
\[
 \alpha_n(L^m)\le\alpha_n(H^m)\alpha_n(K^m).
\]
The three powered positive operators are trace class. For $H^m,K^m$ this follows from $2m>p,q$; for $L^m$ use the displayed ideal bound with $n=1$. Their maximum moduli are their positive real determinant products. Since $r/p+r/q=1$, nonnegative summation yields
\[
 \begin{split}
 M_{L^m}(u)&\le\sum_{n\ge0}\alpha_n(H^m)\alpha_n(K^m)u^n\\
 &\le\left(\sum_{n\ge0}\alpha_n(H^m)u^{nr/p}\right)
       \left(\sum_{n\ge0}\alpha_n(K^m)u^{nr/q}\right).
 \end{split}
\]
The second inequality simply adds all off-diagonal nonnegative terms; it is not a pointwise assertion at a complex argument.

For $\theta,\beta>0$, the substitution $s=u^\theta$ gives $du/u=ds/(\theta s)$ and $u^{-\beta}=s^{-\beta/\theta}$. Therefore
\[
 \int_0^\infty\log M(u^\theta)u^{-1-\beta}du
 =\frac1\theta\int_0^\infty\log M(s)s^{-1-\beta/\theta}ds.
\]
This is valid for nonnegative integrands also when infinite. Here take $\beta=r/(2m)$, and $\theta=r/p$ or $r/q$; the new exponents are $p/(2m)$ and $q/(2m)$, both below one. Writing $c(t)=\pi/\sin(\pi t)$, H0DT.67 gives exactly
\[
 c(r/(2m))S_r(AB)
 \le c(p/(2m))S_p(A)+c(q/(2m))S_q(B).
\]
Indeed $S_{r/(2m)}(L^m)=S_r(AB)$ and $S_{p/(2m)}(H^m)=S_p(A)$, and the Jacobian cancels the reciprocal exponent in the Mellin constant. Divide by $c(r/(2m))$ and let $m\to\infty$: since $c(t)\sim1/t$, the result is $S_r(AB)\le(r/p)S_p(A)+(r/q)S_q(B)$. Rescaling $A$ and $B$ by reciprocal positive constants and minimizing yields the full Schatten Hölder inequality proved independently in H0DT.68. For zero norms it is immediate; otherwise the minimum of $(r/p)t^p S_p(A)+(r/q)t^{-q}S_q(B)$ is $S_p(A)^{r/p}S_q(B)^{r/q}$, because $r/p+r/q=1$. This verifies the entire-function calculation and its use on the following lines; it does not enlarge H0DT to include every unrelated operator assertion on these sheets.
